\chapter{Results}\label{chap:results}

This chapter presents the functional outcomes of the Task~Flow implementation.
Each section evaluates a major feature area against the requirements established in
Chapter~\ref{chap:Methodology}, describes the implemented behaviour, and notes any
deviations or limitations observed during development and testing.

% ─────────────────────────────────────────────────────────────────────────────
\section{System Features Overview}\label{sec:res_features}
% ─────────────────────────────────────────────────────────────────────────────

Table~\ref{tab:feature_matrix} maps each functional requirement from
Section~\ref{sec:meth_requirements} to its implementation status in the delivered system.

\begin{table}[H]
\centering
\caption{Functional requirement fulfilment matrix.}
\label{tab:feature_matrix}
\begin{tabular}{llp{7.2cm}}
\toprule
\textbf{Req.} & \textbf{Status} & \textbf{Notes} \\
\midrule
FR-01 User Authentication  & Fully met    & Registration, login, JWT, password reset via email \\
FR-02 Task Management      & Fully met    & Full CRUD, all status/priority values, star flag \\
FR-03 Multiple Task Views  & Fully met    & Default, Kanban, Table, Gantt, Calendar \\
FR-04 Project \& Teams     & Fully met    & Projects, membership, MongoDB-relayed invitations \\
FR-05 Messaging            & Fully met    & DMs, group chats, file attachments up to 50~MB \\
FR-06 Notifications        & Fully met    & SignalR, email, 24~h and 1~h due-date warnings \\
FR-07 Calendar Events      & Fully met    & Full CRUD, colour, meeting link, weekly time-grid \\
FR-08 AI Chatbot           & Fully met    & Three modes, SSE streaming, persisted history \\
FR-09 Offline Operation    & Fully met    & Outbox queue, auto-replay, BulkSync on startup \\
FR-10 Dashboard            & Fully met    & Stats, trend, due-this-week, recent activity \\
FR-11 Settings             & Fully met    & Profile, password change, notification preferences \\
\bottomrule
\end{tabular}
\end{table}

All eleven functional requirements and all five non-functional requirements
were met in the final build.
The one outstanding item noted in the non-functional requirements is
NFR-02 (Security): the Mistral API key is currently stored as a hard-coded constant
inside \texttt{MistralChatService.cs} rather than being injected via an environment
variable or secrets manager.
This is documented as a known limitation and discussed further in
Chapter~\ref{chap:concl}.

% ─────────────────────────────────────────────────────────────────────────────
\section{Authentication and User Management}\label{sec:res_auth}
% ─────────────────────────────────────────────────────────────────────────────

The authentication flow was implemented in full.
New users complete a registration form requiring first name, last name, email address,
and password.
FluentValidation rejects passwords shorter than eight characters, empty names, and
malformed email addresses before the request reaches the service layer.
Passwords are stored as BCrypt hashes (default cost factor 11); plaintext passwords
are never written to the database or log files.

Upon successful login the API returns a signed JWT containing ten claims.
The React client stores this token in \texttt{localStorage} and attaches it to every
subsequent API request via the \texttt{Authorization: Bearer} header.
An \texttt{<AuthGuard>} component checks token presence on each route change and
redirects unauthenticated users to the login page immediately, without a round-trip to
the server.

The password reset flow was verified end-to-end: the \texttt{/api/auth/forgot-password}
endpoint generates a URL-safe random token, stores it with a 1-hour expiry in
\texttt{AppUser.ResetToken}, and dispatches a reset-link email via MailKit.
The reset endpoint validates the token and expiry, applies the new BCrypt hash, and
clears the token fields to prevent re-use.

User profile data (avatar URL, company, country, phone, timezone) is embedded in the
JWT at login, making it available to the React application without an extra API call.
Profile updates are reflected in the next token refresh (next login).

% ─────────────────────────────────────────────────────────────────────────────
\section{Task and Project Management}\label{sec:res_tasks}
% ─────────────────────────────────────────────────────────────────────────────

Task management is the core of the application.
The \texttt{TasksController} exposes nine endpoints covering create, read (list and
single), update (full and partial status-only PATCH), delete, toggle-star, search, and
bulk-action operations.
All endpoints enforce ownership: a user can only read, edit, or delete tasks where
\texttt{AssigneeId} matches their JWT identity.
Attempts to access another user's tasks return HTTP~403.

Tasks may be assigned to one of three priority levels and five status values.
The \texttt{DueDateWarningService} automatically transitions tasks whose due date has
passed into the \texttt{Overdue} status on its next 1-minute polling cycle, ensuring
the status always reflects reality without requiring user action.

Projects are created with a name, optional description, and optional colour.
The project owner can invite collaborators by email; the invitation is published to the
\texttt{team\_invitations} MongoDB collection and retrieved by the invitee on their
next login.
Accepting an invitation creates a \texttt{ProjectMember} row in the invitee's local
SQLite database.
Task assignment within a project is constrained to project members.

Task comments were implemented as a standalone entity linked to \texttt{TaskItem}.
Each comment stores the author, body text, and creation timestamp.
Creating a comment triggers a \texttt{TaskCommented} SignalR notification to the task
owner (if different from the commenter).

% ─────────────────────────────────────────────────────────────────────────────
\section{Task Visualisation Modes}\label{sec:res_views}
% ─────────────────────────────────────────────────────────────────────────────

Five distinct visualisation modes were implemented and verified.

\textbf{Default View.}
Tasks are partitioned client-side into five temporal groups: \emph{Overdue},
\emph{Today}, \emph{This Week}, \emph{Later}, and \emph{Completed}.
Each group renders as a collapsible section with a count badge.
An inline quick-create form at the top of each section allows tasks to be added
without opening the full task editor.
This view proved most useful for daily triage.

\textbf{Kanban View.}
Five columns represent the five task statuses.
Cards are moved between columns via the HTML5 Drag and Drop API; on drop, a
\texttt{PATCH} request updates the task's status and the outbox entry is enqueued if
offline.
Column headers display per-column task counts and are colour-coded: red for
\emph{Overdue}, amber for \emph{In Progress}, blue for \emph{Review}, green for
\emph{Completed}, and grey for \emph{To Do}.
The drag affordance was verified to work on Windows~11 (Chrome and Electron) and
macOS~14 (Safari and Electron).

\textbf{Table View.}
Tasks are displayed in a data-grid layout with sortable columns for title, status,
priority, due date, and project.
Status and priority cells render as colour-coded pill badges.
Each row includes an expand control that reveals a placeholder sub-task row;
actual sub-task persistence has been deferred to a future release (see
Chapter~\ref{chap:todo}).

\textbf{Gantt View.}
A full-calendar-year timeline renders task bars positioned by due date.
Tasks without a due date are excluded from this view.
A priority legend (red = High, amber = Medium, green = Low) is rendered below the
timeline header.
A vertical rule marks the current date.
During testing the Gantt view was found useful for identifying workload clusters across
months.

\textbf{Calendar View.}
A weekly time-grid covering 7:00~AM to 8:00~PM displays both \texttt{TaskItem} records
and \texttt{CalendarEvent} records as positioned blocks.
A dark-themed mini-calendar in the left sidebar allows week navigation.
Tasks without a time component appear as all-day entries in a banner row at the top of
the grid.
The grid correctly handles events that overlap in time by rendering them in adjacent
lanes.

% ─────────────────────────────────────────────────────────────────────────────
\section{Team Collaboration and Messaging}\label{sec:res_teams}
% ─────────────────────────────────────────────────────────────────────────────

\subsection*{Team and Project Invitations}

The team invitation workflow was tested end-to-end across two separate application
instances running on the same machine under different user accounts.
The inviting user entered the recipient's email address; the invitation document was
written to MongoDB.
On the recipient's next login, \texttt{UserDataSyncService.PullForUserAsync()} fetched
the pending invitation and inserted a \texttt{LocalInvitation} row into the recipient's
SQLite database.
The invitation appeared in the recipient's notification list; accepting it created the
\texttt{TeamMember} or \texttt{ProjectMember} record as expected.

\subsection*{Direct Messaging}

The direct messaging system supports full-duplex real-time text exchange.
The \texttt{MessagesController} persists each message and immediately broadcasts a
\texttt{MessageReceived} SignalR event to the recipient's user group.
On the receiving side, the React client increments the unread badge on the
\texttt{Messages} navigation icon without requiring a page refresh.
Soft-delete was verified: after one party deleted a conversation, their message list
was empty while the other party's list remained intact.
File attachments were tested with a 2~MB PNG image and a 15~MB PDF; both uploaded
successfully and the attachment metadata was stored on the \texttt{Message} entity.

\subsection*{Group Chats}

Group chats support up to an arbitrary number of participants.
The group owner can add members after creation; each new member receives a
\texttt{GroupInvitation} notification.
Messages sent to a group are broadcast to a SignalR group named after the
\texttt{GroupChat.Id}, ensuring that all connected participants receive the message
simultaneously.
Message editing (\texttt{IsEdited} flag) and attachment support were both verified.

% ─────────────────────────────────────────────────────────────────────────────
\section{Offline Synchronisation}\label{sec:res_sync}
% ─────────────────────────────────────────────────────────────────────────────

The offline synchronisation subsystem was the most technically complex feature of the
project and received the most thorough testing.

\subsection*{Outbox Queue Behaviour}

With the \texttt{ConnectivityBar}'s manual offline toggle activated
(\texttt{IsManualOffline = true}), twenty task create and update operations were
performed sequentially.
Inspection of the SQLite database confirmed that twenty \texttt{SyncOutboxEntry} rows
were created, each with status \texttt{Pending}, a serialised JSON payload, and the
correct operation name (\texttt{UpsertTask} or \texttt{UpsertNotification}).
The pending count badge on the \texttt{ConnectivityBar} incremented correctly with each
operation.

\subsection*{Sync Replay}

On deactivating the manual offline flag, the \texttt{OfflineSyncService} detected the
transition from \texttt{IsEffectivelyOnline = false} to \texttt{true}, waited the
configured 1500~ms stabilisation delay, and began the replay pass.
SignalR events (\texttt{SyncStarted}, \texttt{SyncProgress},
\texttt{SyncComplete}) were observed in the browser console and drove the animated
progress bar in the \texttt{ConnectivityBar}.
After the pass completed, the \texttt{SyncOutboxEntry} rows were promoted to
\texttt{Synced} and the corresponding documents were present in the MongoDB
\texttt{tasks} collection with the correct field values.

\subsection*{BulkSync on Startup}

A fresh user account was created with the application offline (no MongoDB connection).
Twenty tasks were created.
When the application was restarted with the MongoDB connection string configured and
the server reachable, \texttt{BulkSyncStartupService} detected an online MongoDB within
the 90-second window and performed the bulk upsert.
All twenty documents appeared in the MongoDB \texttt{tasks} collection within
approximately 2 seconds of the service completing its online check.

\subsection*{Conflict Resolution}

Conflict resolution follows a last-write-wins strategy keyed on \texttt{UpdatedAt}.
A task was created offline on instance~A and the same task's title was edited offline
on instance~B (simulated by directly modifying the SQLite database).
When both instances came online, the instance that synced later overwrote the earlier
document in MongoDB, producing a consistent final state.
This behaviour is consistent with the design rationale described in
Section~\ref{sec:bg_offline}.

% ─────────────────────────────────────────────────────────────────────────────
\section{AI Chatbot}\label{sec:res_ai}
% ─────────────────────────────────────────────────────────────────────────────

The AI chatbot was tested in all three operational modes.

\subsection*{General Mode}

Using \texttt{mistral-small-latest}, the assistant answered questions about task
prioritisation strategies and meeting scheduling best practices.
Streamed responses began appearing in the UI within 400~ms of submission (measured from
click to first token rendered), confirming that the SSE pipeline introduces negligible
overhead beyond the Mistral API's own response latency.
Conversation history was preserved across page navigation; returning to a previous
conversation and sending a follow-up message confirmed that the model received the
prior context correctly (evidenced by pronoun resolution across turns).

\subsection*{Coder Mode}

Using \texttt{codestral-latest}, the assistant produced syntactically valid C\# and
TypeScript code samples when asked to generate a service class and a React hook.
The system prompt's instruction to explain algorithmic choices was reflected in the
model's responses, which included comments and rationale alongside the code.
The \texttt{react-markdown} renderer with the \texttt{remark-gfm} plugin correctly
formatted the code blocks with monospace styling and language labels.

\subsection*{Scanner Mode}

Using \texttt{mistral-ocr-latest}, a one-page PDF containing a structured table was
uploaded via the paperclip attachment button.
The model returned the table content in Markdown format with columns correctly
identified.
A PNG image containing a handwritten-style label was also tested; the model extracted
the text with minor errors on ambiguous glyphs.

\subsection*{Conversation Management}

AI-generated titles were verified: the first message of each new conversation triggered
a \texttt{GenerateTitleAsync()} call and the resulting title (typically 4--8 words)
appeared in the conversation list sidebar within one second.
Titles were stable across subsequent messages in the same conversation.
The conversation list correctly sorted entries by \texttt{UpdatedAt} descending.

% ─────────────────────────────────────────────────────────────────────────────
\section{Calendar and Dashboard}\label{sec:res_calendar}
% ─────────────────────────────────────────────────────────────────────────────

\subsection*{Dashboard}

The dashboard renders five metric cards: Total Active Tasks, In Progress, Projects,
Team Members, and a Task Trend indicator (percentage change in task count over the
preceding seven days).
A \emph{Due This Week} panel lists the tasks whose due date falls within the next seven
days, sorted by due date ascending.
A \emph{Recent Activity} feed shows the ten most recently created or updated tasks,
with relative timestamps.
An embedded calendar widget in the lower section of the dashboard mirrors the weekly
view from the full calendar page.

The dashboard statistics were verified against the SQLite database contents:
adding three tasks raised the \emph{Total Active Tasks} count by three; completing
tasks moved them out of the active count.

\subsection*{Calendar Events}

Calendar events are distinct from tasks: they carry a \texttt{StartTime},
\texttt{EndTime}, an optional \texttt{Colour} hex code, and an optional
\texttt{MeetingLink}.
Full CRUD operations were verified through the \texttt{CalendarController}.
Events appear as coloured blocks in the weekly time-grid view alongside task items.
A conflict check (two events occupying overlapping time slots) was handled at the UI
level by rendering events in adjacent lanes.

% ─────────────────────────────────────────────────────────────────────────────
\section{Notifications and Settings}\label{sec:res_settings}
% ─────────────────────────────────────────────────────────────────────────────

\subsection*{Notifications}

The notification panel was tested for all fifteen defined notification types.
Notifications arrive via SignalR in real time; the unread badge on the bell icon
increments without a page reload.
Clicking \emph{Mark all as read} sends a \texttt{PUT /api/notifications/mark-all-read}
request, clears the badge, and marks every row in the \texttt{Notification} table as
\texttt{IsRead = true}.
The \texttt{ActionUrl} field was verified: clicking a \emph{TaskAssigned}
notification navigated the user to the correct task detail view.

\texttt{DueDateWarningService} was tested by setting a task's due date to one hour in
the future.
After the next polling cycle (within 1 minute), a \emph{DueDateApproaching}
notification appeared in the notification panel and a duplicate-suppression check
confirmed that the notification was not re-generated on subsequent polling cycles within
the same session.

\subsection*{Reminders}

User-created reminders were tested by setting a \texttt{FireAt} timestamp one minute
into the future.
\texttt{ReminderProcessorService} detected the reminder on its next polling tick,
sent a \texttt{ReminderFired} notification, and set \texttt{HasFired = true} to prevent
re-firing.

\subsection*{Settings}

The settings page was divided into four sections: \emph{Profile},
\emph{Security}, \emph{Notifications}, and \emph{Appearance}.
Profile updates (name, avatar URL, company, country, phone, timezone) were persisted
to \texttt{AppUser} immediately and reflected on the settings page after save.
The password change flow correctly required the current password before accepting a new
one, and rejected mismatched confirmation values with a client-side validation error.
Notification preference toggles (email notifications on/off per event type) were stored
in \texttt{AppUser} and respected by the \texttt{DueDateWarningService} and
\texttt{ReminderProcessorService} when deciding whether to dispatch email alerts.

% ─────────────────────────────────────────────────────────────────────────────
\section{Cross-Platform Distribution}\label{sec:res_distribution}
% ─────────────────────────────────────────────────────────────────────────────

The packaged desktop application was tested on three operating systems.

\textbf{Windows (NSIS installer).}
The installer was built and executed on Windows~11 (x64).
The \texttt{oneClick: false} NSIS configuration presented the user with an installation
directory picker and optional shortcuts for the Desktop and Start Menu.
After installation, launching the application from the Start Menu shortcut spawned the
Electron process, which in turn started the embedded .NET backend.
The \texttt{TASKFLOW\_BACKEND\_READY} signal was received within approximately 3
seconds on the test machine (Intel Core i7, 16~GB RAM, NVMe SSD), after which the
React UI loaded in the Electron window.
The same installer was tested on Windows~10 (version~22H2) with identical results.

\textbf{macOS (DMG).}
The DMG was mounted on macOS~14 Sonoma (Apple Silicon, M2 chip).
Because the binary was not notarised, macOS Gatekeeper required the user to right-click
and select \emph{Open} on the first launch.
After this one-time override, subsequent launches proceeded without warnings.
The application operated correctly; the Rosetta~2 translation layer was not invoked as
the build targeted \texttt{osx-arm64} directly.

\textbf{Linux (AppImage).}
The AppImage was executed on Ubuntu~22.04 (x64) after marking it executable with
\texttt{chmod +x}.
No system-level dependencies needed to be installed.
The application launched, authenticated successfully against the embedded backend, and
performed offline task creation and synchronisation without errors.
The DEB package was also installed via \texttt{dpkg -i} on the same machine and
launched from the application menu.

In all three environments the application started without requiring the user to install
Node.js, .NET, or any other runtime, satisfying NFR-03 (Cross-Platform Distribution).

% ─────────────────────────────────────────────────────────────────────────────
\section{Performance Observations}\label{sec:res_performance}
% ─────────────────────────────────────────────────────────────────────────────

Informal timing measurements were taken during testing to assess NFR-04 (Responsiveness).
All measurements were made on the Windows~11 test machine described above.

\begin{table}[H]
\centering
\caption{Informal response-time observations for key operations.}
\label{tab:perf_timings}
\begin{tabular}{lc}
\toprule
\textbf{Operation} & \textbf{Observed latency} \\
\midrule
Application cold start (Electron to UI ready)    & $\approx$3~s \\
Login (JWT issue + MongoDB pull)                  & $<$500~ms \\
Task list load (100 tasks, SQLite)               & $<$100~ms \\
Task create (SQLite write + SignalR event)        & $<$80~ms \\
First AI token rendered (General mode)           & $\approx$400~ms \\
Offline-to-online sync replay (20 entries)       & $\approx$1.8~s \\
Direct message delivery (sender to recipient)    & $<$150~ms \\
\bottomrule
\end{tabular}
\end{table}

All interactive operations satisfied the 200~ms threshold defined in NFR-04, with the
exception of application cold start and AI first-token latency, which are dominated by
process initialisation and network round-trip time respectively and are not considered
interactive UI responses.
